\documentclass[12pt,a4paper]{article}
\usepackage{fontspec}
\setmainfont{DejaVu Serif}
\setsansfont{DejaVu Sans}
\setmonofont{DejaVu Sans Mono}
\usepackage{geometry}
\usepackage{hyperref}
\usepackage{listings}
\usepackage{xcolor}
\usepackage{booktabs}
\usepackage{graphicx}
\usepackage{parskip}
\usepackage{enumitem}
\usepackage{amsmath}
\usepackage{tcolorbox}
\usepackage{fancyhdr}
\usepackage{titlesec}
\usepackage{array}
\usepackage{ragged2e}
\usepackage{setspace}

\geometry{a4paper, left=25mm, right=25mm, top=32mm, bottom=30mm}

\setlength{\headheight}{14pt}

\pagestyle{fancy}
\fancyhf{}
\rhead{\small CodeAlchemist -- 6. Hafta Raporu}
\lhead{\small Dila KEMER | Bilgisayar Mühendisliği}
\cfoot{\thepage}
\renewcommand{\headrulewidth}{0.4pt}

\definecolor{codegray}{rgb}{0.5,0.5,0.5}
\definecolor{backcolour}{rgb}{0.95,0.95,0.92}
\definecolor{keywordblue}{rgb}{0.0,0.2,0.7}
\definecolor{commentgreen}{rgb}{0.1,0.5,0.1}
\definecolor{stringpurple}{rgb}{0.5,0.0,0.5}
\definecolor{noteblue}{rgb}{0.88,0.93,1.0}
\definecolor{noteborder}{rgb}{0.3,0.5,0.8}
\definecolor{ozetsarisi}{rgb}{0.98,0.96,0.88}
\definecolor{ozetborder}{rgb}{0.7,0.5,0.1}

\titleformat{\section}{\large\bfseries\color{keywordblue}}{\thesection}{1em}{}[\titlerule]
\titleformat{\subsection}{\normalsize\bfseries}{\thesubsection}{1em}{}

\newcolumntype{P}[1]{>{\RaggedRight\arraybackslash}p{#1}}

\lstdefinestyle{mystyle}{
    backgroundcolor=\color{backcolour},
    commentstyle=\color{commentgreen},
    keywordstyle=\color{keywordblue}\bfseries,
    numberstyle=\tiny\color{codegray},
    stringstyle=\color{stringpurple},
    basicstyle=\ttfamily\footnotesize,
    breaklines=true,
    captionpos=b,
    keepspaces=true,
    numbers=left,
    numbersep=5pt,
    showspaces=false,
    tabsize=2,
    frame=single
}
\lstset{style=mystyle}

\lstdefinelanguage{JavaScript}{
  keywords={break,case,catch,continue,debugger,default,delete,do,else,finally,
    for,function,if,in,instanceof,new,return,switch,this,throw,try,typeof,
    var,void,while,with,const,let,async,await,of,class,import,export,from},
  morestring=[b]",
  morestring=[b]',
  morestring=[b]`,
  morecomment=[l]{//},
  morecomment=[s]{/*}{*/},
  sensitive=true
}

\tcbuselibrary{skins,breakable}
\newtcolorbox{infobox}[1][]{
  colback=noteblue,
  colframe=noteborder,
  fonttitle=\bfseries\small,
  title={#1},
  breakable
}

\newtcolorbox{ozetbox}[1][]{
  colback=ozetsarisi,
  colframe=ozetborder,
  fonttitle=\bfseries\normalsize,
  title={#1},
  breakable,
  boxrule=1pt,
  arc=4pt
}

\hypersetup{
  colorlinks=true,
  linkcolor=keywordblue,
  urlcolor=keywordblue,
  pdftitle={CodeAlchemist 6. Hafta Raporu},
  pdfauthor={Dila KEMER}
}

% ============================================================
\begin{document}
\sloppy

%------------------------------------------------------------
% KAPAK SAYFASI
%------------------------------------------------------------
\begin{titlepage}
  \centering
  \vspace*{1.8cm}

  {\Huge\bfseries\color{keywordblue} CodeAlchemist}\\[0.8em]
  {\large\bfseries Yapay Zeka Destekli Kod Geliştirme Platformu}\\[2em]

  \rule{\textwidth}{1pt}\\[0.8em]
  {\LARGE\bfseries 6. ve 7. Hafta -- Birleştirilmiş Teknik Rapor}\\[0.5em]
  \rule{\textwidth}{1pt}\\[2em]

  {\large Haftalık Özellik Geliştirme, Mimari Kararlar ve Etki Analizi}

  \vspace{2.5em}

  \begin{tabular}{rl}
    \textbf{Hazırlayan:}   & Dila KEMER \\[0.4em]
    \textbf{Üniversite:}   & İzmir Bakırçay Üniversitesi \\[0.4em]
    \textbf{Bölüm:}        & Bilgisayar Mühendisliği \\[0.4em]
    \textbf{Sınıf:}        & 4. Sınıf \\[0.4em]
    \textbf{Ders:}         & İşletmede Mesleki Eğitim (İME) \\[0.4em]
    \textbf{Rapor Türü:}   & 2 Haftalık Teknik Konsolidasyon \\[0.4em]
    \textbf{Tarih:}        & Mart 2026 \\
  \end{tabular}

  \vspace{2.4em}

  \begin{ozetbox}[Canlı Demo]
    Projenin çalışır halini aşağıdaki adresten inceleyebilirsiniz:\\[0.6em]
    \begin{center}
      \large\href{https://code-alchemist-last-version.onrender.com/}%
      {https://code-alchemist-last-version.onrender.com/}
    \end{center}
  \end{ozetbox}

  \vfill
  {\small\color{gray} CodeAlchemist AI Projesi -- Geliştirici Ekip}
\end{titlepage}

\thispagestyle{empty}

\newpage
\section*{Özet}
\addcontentsline{toc}{section}{Özet}

Bu döküman, CodeAlchemist projesinin \textbf{6. hafta} ve \textbf{7. hafta} geliştirmelerini birbirinden ayrıştırılmış iki teknik rapor halinde sunar. 6. hafta daha çok geliştirici üretkenliği, bağlam görünürlüğü ve ölçümleme odaklı altyapı özelliklerini içerirken; 7. hafta kullanıcı bağlılığı, sosyal etkileşim, görselleştirme ve kişiselleştirme başlıklarında ürün deneyimini güçlendiren özelliklere odaklanmıştır.

Rapor; her hafta için kapsam, mimari, bileşen sorumlulukları, API akışları, veri modeli etkisi, test/validasyon çıktıları ve bilinen iyileştirme alanlarını detaylandırır.

\textbf{Anahtar Kelimeler:} Apply Patch, Context Window, TTFT, Cost Dashboard, Gamification, Weekly Report, Collaboration, Theme Store.

\section*{Abstract}
\addcontentsline{toc}{section}{Abstract}

This document delivers a split technical report for the 6th and 7th development weeks of the CodeAlchemist platform. Week 6 focuses on developer productivity and observability features; Week 7 focuses on engagement, social sharing, personalization, and analytics UX.

Each weekly part includes scope, architecture, component responsibilities, API behavior, data model impact, validation outcomes, and future improvements.

\textbf{Keywords:} Apply Patch, Context Window, TTFT, Cost Dashboard, Gamification, Weekly Report, Collaboration, Theme Store.

\newpage
\tableofcontents
\newpage

%====================================================================
\section{Proje Bağlamı ve Teknik Zemin}
%====================================================================

CodeAlchemist; React tabanlı modern bir frontend, Flask tabanlı API katmanı, SQLAlchemy ORM ve çoklu model yönlendirme bileşenlerinden oluşan bir AI destekli kod asistan platformudur.

\textbf{Teknoloji Özeti:}
\begin{itemize}[leftmargin=*]
  \item \textbf{Frontend:} React 18, Vite, TailwindCSS
  \item \textbf{Backend:} Flask, SQLAlchemy, JWT tabanlı auth
  \item \textbf{Veri Katmanı:} SQLite (geliştirme), ilişkisel modelleme
  \item \textbf{AI Katmanı:} Gemini, GPT, Claude modelleri için yönlendirme altyapısı
  \item \textbf{Gerçek zamanlılık:} SSE tabanlı token akışı ve canlı sohbet güncellemeleri
\end{itemize}

\begin{infobox}[Rapor Yapısı]
Bu belge iki ana rapora bölünmüştür:
\textbf{Bölüm I: 6. Hafta Eklenen Özellikler} ve
\textbf{Bölüm II: 7. Hafta Eklenen Özellikler}.
Her bölüm kendi içinde bağımsız okunabilecek şekilde tasarlanmıştır.
\end{infobox}

\newpage
%====================================================================
\section{Bölüm I -- 6. Hafta Eklenen Özellikler (Detaylı Rapor)}
%====================================================================

\subsection{Hafta 6 Kapsam Özeti}

6. haftadaki geliştirmeler, doğrudan geliştirici iş akışını hızlandıran ve AI etkileşim kalitesini görünür kılan özellikler üzerine kurulmuştur:
\begin{enumerate}[leftmargin=*, label=\textbf{H6.\arabic*.}]
  \item Apply Patch (otonom kod uygulama + undo)
  \item Context Bar (bağlam/token yükü görünürlüğü)
  \item Yanıt Süresi İndikatörü (response time + TTFT)
  \item Model Maliyet Paneli (kullanım ve tahmini maliyet)
  \item Karşılama Sayfası ve kimlik doğrulama akışı iyileştirmeleri
\end{enumerate}

\subsection{H6.1 -- Apply Patch (Otonom Kod Uygulama)}

\textbf{Hedef:} AI tarafından önerilen kodun editöre kontrollü şekilde uygulanması ve geri alınabilirliğin korunması.

\textbf{Mimari Karar:}
\begin{itemize}[leftmargin=*]
  \item Uygulama noktası frontendde tutuldu, böylece anlık geri bildirim sağlandı.
  \item Undo için snapshot tabanlı LIFO yaklaşımı benimsendi.
  \item Kısmi kabul için görsel diff tabanlı modal akışı korundu.
\end{itemize}

\textbf{Kritik Bileşenler:}
\begin{itemize}[leftmargin=*]
  \item \texttt{client/src/components/ChatInterface.jsx}
  \item \texttt{client/src/components/InteractiveMergeModal.jsx}
\end{itemize}

\begin{lstlisting}[language=JavaScript, caption=Apply Patch akışında kullanılan temel state yapıları]
const [patchHistory, setPatchHistory] = useState([]);
const [patchModal, setPatchModal] = useState({ show: false, newCode: '', type: 'replace' });
const [mergeModal, setMergeModal] = useState({ isOpen: false, newCode: '' });
\end{lstlisting}

\textbf{Etkisi:} Kopyala-yapıştır hatalarını azaltır, kullanıcıya güvenli deneme alanı verir.

\subsection{H6.2 -- Context Bar (Bağlam Görünürlüğü)}

\textbf{Hedef:} Prompt + code + ek bağlam kaynaklarının model sınırına ne kadar yaklaştığını canlı göstermek.

Token tahmini formülü:
$$\text{tokenCount} = \left\lceil\frac{|question| + |code|}{4}\right\rceil$$

\textbf{Renk Eşikleri:}
\begin{itemize}[leftmargin=*]
  \item $<50\%$: güvenli kullanım
  \item $50\%-80\%$: dikkat seviyesi
  \item $>80\%$: kritik bölge
\end{itemize}

\textbf{Etkisi:} Kullanıcı, modelin bağlam limitine yaklaşmadan girdiyi optimize edebilir.

\subsection{H6.3 -- AI Yanıt Süresi ve TTFT Ölçümü}

\textbf{Hedef:} Model performansını kullanıcıya şeffaf şekilde göstermek.

\begin{lstlisting}[language=JavaScript, caption=Toplam süre ve TTFT ölçüm prensibi]
const startTime = Date.now();
let firstChunkTime = null;
// ... SSE stream
const totalDuration = (endTime - startTime) / 1000;
const ttft = firstChunkTime ? (firstChunkTime - startTime) / 1000 : null;
\end{lstlisting}

\textbf{Etkisi:} Model seçimi ve altyapı kıyaslaması için görünür metrik sağlar.

\subsection{H6.4 -- Model Maliyet Paneli}

\textbf{Hedef:} Kullanılan model dağılımını ve tahmini maliyeti anlamlandırmak.

\textbf{Maliyet Yaklaşımı:}
$$
\text{cost} = \frac{N \cdot T_{in} \cdot P_{in} + N \cdot T_{out} \cdot P_{out}}{1{,}000{,}000}
$$

\textbf{Bileşenler:}
\begin{itemize}[leftmargin=*]
  \item \texttt{client/src/components/ModelCostDashboard.jsx}
  \item \texttt{server/app.py} (istatistik endpointleri)
\end{itemize}

\textbf{Etkisi:} Ücretsiz/ücretli model stratejileri için operasyonel içgörü üretir.

\subsection{H6.5 -- Karşılama Sayfası ve Auth Akışı}

\textbf{Hedef:} Oturumsuz kullanıcı için ürün değerini hızlı anlatan modern bir giriş yüzeyi sağlamak.

\textbf{Kritik Noktalar:}
\begin{itemize}[leftmargin=*]
  \item Landing page görünürlüğü localStorage oturum durumuyla yönetildi.
  \item Auth modal ve landing katmanları z-index ile çakışmasız çalışacak şekilde ayrıldı.
  \item Girişten vazgeçen kullanıcıya tekrar deneme akışı bırakıldı.
\end{itemize}

\subsection{Hafta 6 Sonuç ve Teknik Değerlendirme}

6. hafta çıktıları, geliştirici deneyiminde iki temel kazanım üretmiştir:
\begin{itemize}[leftmargin=*]
  \item \textbf{Operasyonel hız:} Apply Patch + Context Bar kombinasyonu ile daha hızlı ve kontrollü iterasyon.
  \item \textbf{Gözlemlenebilirlik:} TTFT ve maliyet paneli ile ölçülebilir kullanım davranışı.
\end{itemize}

\newpage
%====================================================================
\section{Bölüm II -- 7. Hafta Eklenen Özellikler (Detaylı Rapor)}
%====================================================================

\subsection{Hafta 7 Kapsam Özeti}

7. hafta geliştirmeleri, doğrudan kullanıcı bağlılığı ve ürün içi etkileşimleri güçlendirmeye odaklanmıştır:
\begin{enumerate}[leftmargin=*, label=\textbf{H7.\arabic*.}]
  \item Gamification paneli ve XP/level doğruluk iyileştirmeleri
  \item Haftalık rapor ekranı (bar/line görünüm, okunabilirlik artışı)
  \item Collaboration paylaşım seçenekleri (link, open, WhatsApp, e-posta)
  \item Theme Store güvenliği ve oturum-temelli tema davranışı düzeltmeleri
\end{enumerate}

\subsection{H7.1 -- Gamification: Seviye, Rank ve İlerleme Doğruluğu}

\textbf{Problem:} Bazı ekranlarda level değeri veritabanındaki stale alanlardan okunarak XP ile tutarsız gösteriliyordu.

\textbf{Çözüm Stratejisi:}
\begin{itemize}[leftmargin=*]
  \item Frontendde level değerini doğrudan XP'den türeten yaklaşım benimsendi.
  \item Backendde profil endpointinde seviye alanı XP bazında normalize edildi.
  \item Kullanıcı kimliği, username ve rank title gösterimi standartlaştırıldı.
\end{itemize}

\textbf{Matematiksel Model:}
$$
\text{level} = \left\lfloor\sqrt{\frac{xp}{100}}\right\rfloor + 1
$$

\textbf{Etki:} ``Seviye 6 ama progress 0\%'' gibi çelişkili durumlar kaldırıldı.

\subsection{H7.2 -- Weekly Report: Grafik Okunabilirlik ve Mod Değişimi}

\textbf{Yapılan İyileştirmeler:}
\begin{itemize}[leftmargin=*]
  \item Çubuk grafikte her gün için görünür konteyner ve minimum yükseklik kuralları tanımlandı.
  \item Çizgi modunda karmaşık katmanlar sadeleştirilerek temiz SVG polyline yapısına geçildi.
  \item Nokta/marker gürültüsü azaltılarak çizgi okunabilirliği yükseltildi.
  \item Min/Max başlıkları ve gün etiketleri netleştirildi.
\end{itemize}

\textbf{Bileşen:} \texttt{client/src/components/WeeklyReport.jsx}

\subsection{H7.3 -- Collaboration Share: Çok Kanallı Paylaşım}

\textbf{Hedef:} Üretilen işbirliği bağlantısının farklı kanallardan hızlı paylaşımı.

\textbf{Eklenen Opsiyonlar:}
\begin{itemize}[leftmargin=*]
  \item Link kopyalama
  \item Yeni sekmede açma
  \item WhatsApp paylaşımı
  \item E-posta paylaşımı
\end{itemize}

\textbf{Teknik Not:} URL parametreleri \texttt{encodeURIComponent} ile güvenli biçimde kodlandı.

\textbf{Bileşen Akışı:}
\begin{itemize}[leftmargin=*]
  \item \texttt{client/src/App.jsx} içinde paylaşım linkinin üretilmesi
  \item Aynı dosyada aksiyon modalı ile kanal seçimi
\end{itemize}

\subsection{H7.4 -- Theme Store: Oturum Güvenliği ve Sahiplik Davranışı}

\textbf{Problem:} Kullanıcı çıkış yapsa bile aktif premium tema veya ``Applied'' durumu arayüzde kalabiliyordu.

\textbf{Alınan Önlemler:}
\begin{itemize}[leftmargin=*]
  \item Logout sırasında frontendde token/user/theme temizliği zorunlu hale getirildi.
  \item DOM üzerindeki \texttt{data-theme} değeri logout anında varsayılan temaya çekildi.
  \item Theme Store işlemlerinde token yoksa apply/unlock bloklandı.
  \item Logout öncesi backendde aktif tema varsayılan temaya set edilerek kalıcı tutarlılık sağlandı.
\end{itemize}

\textbf{İlgili Dosyalar:}
\begin{itemize}[leftmargin=*]
  \item \texttt{client/src/App.jsx}
  \item \texttt{client/src/components/ThemeStore.jsx}
\end{itemize}

\subsection{H7.5 -- UI Dokunuşları ve Kullanılabilirlik}

Bu hafta boyunca küçük ama etkili UI düzenlemeleri de yapıldı:
\begin{itemize}[leftmargin=*]
  \item ``Magical Tools'' butonunda ikon görünürlüğü iyileştirildi.
  \item Haftalık rapor çizgi/çubuk geçişi görsel olarak daha okunur hale getirildi.
  \item Gamification kartlarında kimlik ve rank bilgileri daha anlaşılır etiketlendi.
\end{itemize}

\subsection{Hafta 7 Sonuç ve Teknik Değerlendirme}

7. hafta çıktıları üç ana başlıkta değer üretmiştir:
\begin{itemize}[leftmargin=*]
  \item \textbf{Tutarlılık:} XP, level, rank ve tema sahipliği verilerinde ekranlar arası uyum.
  \item \textbf{Paylaşılabilirlik:} Collaboration linklerinin farklı kanallara taşınması.
  \item \textbf{Okunabilirlik:} Haftalık metriklerin daha doğru ve anlaşılır görselleştirilmesi.
\end{itemize}

\newpage
%====================================================================
\section{Hafta 6 ve Hafta 7 Karşılaştırmalı Analiz}
%====================================================================

\begin{table}[h!]
  \centering\small
  \begin{tabular}{@{} P{2.8cm} P{4.9cm} P{4.9cm} @{}}
    \toprule
    \textbf{Kıyas} & \textbf{6. Hafta} & \textbf{7. Hafta} \\
    \midrule
    Odak & Geliştirici üretkenliği, ölçümleme & Kullanıcı bağlılığı, paylaşım, kişiselleştirme \\
    Ağır Katman & Frontend (UX + state) & Fullstack (UI + API + veri tutarlılığı) \\
    Çıktı Türü & Operasyonel araçlar & Ürün davranışı ve topluluk etkileşimi \\
    Kritik Risk & Token tahmin doğruluğu & Oturumsuz tema/kimlik tutarsızlığı \\
    Ölçülebilir Kazanım & Patch hızlanması, performans görünürlüğü & Daha net grafikler, daha güvenli kullanıcı deneyimi \\
    \bottomrule
  \end{tabular}
  \caption{6. ve 7. hafta teknik odak farkları}
\end{table}

\section{Sonuç ve Sonraki Adımlar}

Bu iki haftalık geliştirme periyodu sonunda CodeAlchemist, hem mühendislik verimliliği hem de son kullanıcı deneyimi tarafında belirgin olgunlaşma kazanmıştır. Özellikle 7. haftada yapılan tutarlılık düzeltmeleri (XP/level, tema sahipliği, paylaşım akışları), ürünün güvenilirlik algısını doğrudan artırmıştır.

\textbf{Önerilen kısa vadeli yol haritası:}
\begin{itemize}[leftmargin=*]
  \item Theme sahipliği doğrulamasının backend kaynaklı tek bir ``source of truth'' endpointi ile standartlaştırılması
  \item Weekly report için 0-aktivite günlerinin dim bar ile gösterilmesi
  \item Gamification için başarı rozetlerinin kazanım tetikleyicilerinin test kapsamına alınması
  \item Collaboration paylaşımı için kullanım analitikleri (hangi kanal daha çok kullanılıyor) eklenmesi
\end{itemize}

\begin{ozetbox}[Canlı Erişim]
\begin{center}
\large
\href{https://code-alchemist-last-version.onrender.com/}{https://code-alchemist-last-version.onrender.com/}
\end{center}
\end{ozetbox}

\vspace{1em}
\noindent{\small\color{gray}
Referanslar: \texttt{client/src/App.jsx}, \texttt{client/src/components/WeeklyReport.jsx},
\texttt{client/src/components/GamificationPanel.jsx}, \texttt{client/src/components/ThemeStore.jsx},
\texttt{client/src/components/ChatInterface.jsx}, \texttt{server/app.py}.}

\end{document}